%% toc: P15 164 164.12
\section{Two source comparisons: derivations}
\label{sec:p15-w02b-p15-025-annex}

Throughout, \(r\) is a range measure in the sense of
\cref{def:p15-power-range-measure}, \(g=1-r\), and \(P_r\) is the
probability of \cref{def:p15-w02b-p15-025-common-model}. The laws of
\cref{thm:p15-power-range-laws} are used without further mention:
\(r(\bot)=0\), \(r(\neg U)=1-r(U)\), equivalent sentences have equal
values, entailment gives \(r(U)\le r(V)\), and
\(r(U\vee V)=r(U)+r(V)-r(U\wedge V)\).

\begin{annexitem}{annex:p15-w02b-p15-025-common-model}{thm:p15-w02b-p15-025-common-model}{The range measure as a probability}
\begin{proof}
If \([U]=[V]\) then \(U\equiv V\) and \(r(U)=r(V)\), so \(P_r\) is well
defined. The class of valid sentences is the top element, and \(r\) is
\(1\) there. Two classes \([U],[V]\) are disjoint in \(\mathcal B\)
exactly when \(U\wedge V\) is contradictory; then
\(P_r([U]\vee[V])=r(U\vee V)=r(U)+r(V)\) by the additivity clause of
\cref{def:p15-power-range-measure}. So \(P_r\) is a finitely additive
probability on \(\mathcal B\).

If \(r(S)>0\), then \(p_r(T,S)=r(T\wedge S)/r(S)=P_r(T\mid S)\) by the
definition of \(p_r\) and of conditional probability. If \(r(S)<1\),
then \(g(S)=r(\neg S)>0\), and \cref{thm:p15-power-probability-dual}
gives \(s_r(T,S)=p_r(\neg T,\neg S)=r(\neg T\wedge\neg S)/r(\neg S)
=P_r(\neg T\mid\neg S)\). Finally \(p_r(T,S)\) is defined exactly when
\(r(S)>0\), and \(s_r(T,S)\) exactly when \(g(S)>0\), that is,
\(r(S)<1\) (\cref{def:p15-power-dual-measures}).
\end{proof}
\end{annexitem}

\begin{annexitem}{annex:p15-w02b-p15-025-prediction}{thm:p15-w02b-p15-025-prediction}{Positive power and prediction}
\begin{proof}
Write \(\Omega\), \(D_v\), \(M_v\), \([U]\) and \(\mathcal C(U)\) as in
\cref{def:p15-power-language}. Two facts are used.

(a) \emph{Every ground assignment is realized.} For \(v\in\Omega\) let
\(\mathfrak A_v\) have the constants of the language as its domain (a single
new element if there are none), each constant naming itself, and let a ground atom be true in
\(\mathfrak A_v\) exactly when \(v\) gives it value~\(1\). The language
has no identity, so this is a structure, and a singular sentence is true
in it exactly when it is true under \(v\).

(b) \emph{If \(U\) is singular and \(v\notin[U]\), then \(U\entails M_v\).}
By \cref{thm:p15-power-normal-forms}, \(U\) is equivalent to the
conjunction of \(\mathcal C(U)=\{M_u:u\notin[U]\}\), which contains
\(M_v\).

Item~1, from right to left. Suppose \(s_0(T,K)>0\). By
\cref{thm:p15-power-optimal-split}, \(s_0(T,K)=|K_T|/|\mathcal C(S)|\),
so there is \(M_v\in\mathcal C(S)\) with \(T\entails M_v\). Put
\(E_2=\{M_v\}\) and \(E_1=\mathcal C(S)\setminus\{M_v\}\). Then
\(E_1\cup\{T\}\entails M_v\). Each \(M_u\) with \(u\neq v\) is false
only under \(u\), so every member of \(E_1\) is true under \(v\), while
\(M_v=\neg D_v\) is false under \(v\). In \(\mathfrak A_v\) the members
of \(E_1\) are true and \(M_v\) is false, so \(E_1\not\entails M_v\).
(When \(E_1\) is empty this says that \(M_v\) is not valid.) The report
\(\mathcal C(S)\) is therefore predictive for \(T\).

Item~1, from left to right. Let \(\mathcal C(S)=E_1\mathbin{\dot\cup}E_2\)
with \(E_2\) nonempty and every member of \(E_2\) derivable from
\(E_1\cup\{T\}\). The members of \(\mathcal C(S)\) are distinct minimal
sentences, so \cref{thm:p15-power-optimal-split} (source
Theorem~8.2) shows that every member of \(E_2\) follows from \(T\)
alone. Hence \(K_T\supseteq E_2\neq\varnothing\) and \(s_0(T,K)>0\).
This direction does not use the clause ``not from \(E_1\) alone''.

Item~2. Let \(K=E_1\mathbin{\dot\cup}E_2\) witness predictivity and pick
\(e\in E_2\). Let \(S_1\) be the conjunction of \(E_1\) (\(\top\) when
\(E_1\) is empty). From \(E_1\cup\{T\}\entails e\) we get
\(T\entails\neg S_1\vee e\), a singular sentence. From
\(E_1\not\entails e\) there is a structure in which \(S_1\) is true and
\(e\) false; let \(v\) be its ground assignment. Then \(v\) falsifies
\(\neg S_1\vee e\), so by (b) \(T\entails\neg S_1\vee e\entails M_v\).
Since \(e\in K\) is false under \(v\), so is \(S\); hence \(v\notin[S]\)
and \(M_v\in\mathcal C(S)\). Thus \(K_T\neq\varnothing\) and
\(s_0(T,K)>0\) as before.

Item~3. Here \(T=\forall x(P(x)\to Q(x))\) and
\[
 K_2=\{P(a)\wedge Q(a),\ \neg P(b)\wedge\neg Q(b),\ Q(c),\ \neg P(d)\},
\]
and \(s_0(T,K_2)=25/36\) by \cref{thm:p15-power-source-examples}. For
each member we give a structure with domain \(\{a,b,c,d\}\), each
constant naming itself, in which \(T\) and the other three members are
true and that member is false. In each case \(P\subseteq Q\), so \(T\)
holds.
\begin{itemize}
\item \(P=\varnothing\), \(Q=\{c\}\): \(P(a)\wedge Q(a)\) is false;
the others are true.
\item \(P=\{a\}\), \(Q=\{a,b,c\}\): \(\neg P(b)\wedge\neg Q(b)\) is
false; the others are true.
\item \(P=\{a\}\), \(Q=\{a\}\): \(Q(c)\) is false; the others are true.
\item \(P=\{a,d\}\), \(Q=\{a,c,d\}\): \(\neg P(d)\) is false; the others
are true.
\end{itemize}
So no member of \(K_2\) follows from the remaining members together with
\(T\). If \(K_2=E_1\mathbin{\dot\cup}E_2\) witnessed predictivity, any
\(e\in E_2\) would follow from \(E_1\cup\{T\}\subseteq
(K_2\setminus\{e\})\cup\{T\}\), which has just been excluded. Hence
\(K_2\) is not predictive for \(T\), although \(s_0(T,K_2)>0\).
\end{proof}
\end{annexitem}

\begin{annexitem}{annex:p15-w02b-p15-025-increment}{thm:p15-w02b-p15-025-increment}{Power above content}
\begin{proof}
Let \(0<r(S)<1\). First,
\[
 \Delta_r(\neg T,\neg S)
 =r(\neg T\wedge\neg S)-(1-r(T))(1-r(S)).
\]
Since \(\neg T\wedge\neg S\equiv\neg(T\vee S)\),
\(r(\neg T\wedge\neg S)=1-r(T)-r(S)+r(T\wedge S)\). Expanding the
product gives \(\Delta_r(\neg T,\neg S)=r(T\wedge S)-r(T)r(S)
=\Delta_r(S,T)\). By \cref{thm:p15-w02b-p15-025-common-model},
\[
 \begin{aligned}
 s_r(T,S)-g(T)&=\frac{r(\neg T\wedge\neg S)}{r(\neg S)}-r(\neg T)\\
 &=\frac{r(\neg T\wedge\neg S)-r(\neg T)r(\neg S)}{r(\neg S)}
 =\frac{\Delta_r(S,T)}{r(\neg S)} .
 \end{aligned}
\]
The middle expression is \(P_r(\neg T\mid\neg S)-P_r(\neg T)\), which
is \(R_D(\neg S\Yright\neg T)_r\) by \cref{def:p15-r43-measures}, since
\(P_r(\neg S)>0\). Also
\(R_D(S\Yright T)_r=r(T\wedge S)/r(S)-r(T)=\Delta_r(S,T)/r(S)\), so
\[
 \frac{\Delta_r(S,T)}{r(\neg S)}
 =\frac{r(S)}{1-r(S)}\cdot\frac{\Delta_r(S,T)}{r(S)}
 =O\bigl(r(S)\bigr)\,R_D(S\Yright T)_r .
\]

Item~1. \(r(\neg S)>0\), so \(s_r(T,S)-g(T)\) has the sign of
\(\Delta_r(S,T)\), which by \cref{def:p15-r43-polarity} is positive,
zero or negative exactly when \(S\) is positive, irrelevant or negative
to \(T\) under \(P_r\).

Item~2. In \cref{def:p15-power-dual-syntax}, with \(T,S\) as sentence
metavariables, \(D(T)=T\), \(D(S)=S\) and \(D(\top)=\bot\). Replacing
each \(p_r(B,C)\) by \(s_r(D(B),D(C))\) sends
\(p_r(T,S)-p_r(T,\top)\) to \(s_r(T,S)-s_r(T,\bot)\); numerical
subtraction is unchanged. Further \(g(\bot)=1-r(\bot)=1\) and
\(T\vee\bot\equiv T\), so \(s_r(T,\bot)=g(T)/1=g(T)\). The term
\(p_r(T,\top)=r(T\wedge\top)/r(\top)=r(T)\), so
\(p_r(T,S)-p_r(T,\top)=\Delta_r(S,T)/r(S)\). Both differences are
therefore \(\Delta_r(S,T)\) divided by a positive number.

Item~3. If \(r(S)=0\), then \(g(S)=1\) and \(r(T\wedge S)\le r(S)=0\).
Hence \(g(T\vee S)=1-r(T)-r(S)+r(T\wedge S)=1-r(T)=g(T)\), and
\(s_r(T,S)=g(T)/1=g(T)\).

Item~4. Let \(r\) give each of the eight assignments to \(A,B,C\) the
weight \(1/8\). Then \(r(\neg A\wedge\neg B)=1/4\), \(r(\neg B)=1/2\),
and \(s_r(A,B)=(1/4)/(1/2)=1/2\). Also
\(\neg(A\wedge C)\wedge\neg B\) holds at three of the four assignments
with \(B\) false, so \(s_r(A\wedge C,B)=(3/8)/(1/2)=3/4\). Further
\(\Delta_r(B,A)=1/4-1/4=0\) and \(\Delta_r(B,A\wedge C)=1/8-1/8=0\), so
\(B\) is irrelevant to both targets. If some \(\iota\) represented
quality, irrelevance would force \(s_r(A,B)=\iota=s_r(A\wedge C,B)\),
that is, \(1/2=3/4\). No such \(\iota\) exists.
\end{proof}
\end{annexitem}

\begin{annexitem}{annex:p15-w02b-p15-025-conclusive}{thm:p15-w02b-p15-025-conclusive}{Power one and conclusive relevance}
\begin{proof}
Let \(r(T)>0\) and \(0<r(S)<1\).

(1)\(\Leftrightarrow\)(2). By \cref{thm:p15-w02b-p15-025-common-model},
\(s_r(T,S)=r(\neg T\wedge\neg S)/r(\neg S)\), and
\(r(\neg S)=r(\neg T\wedge\neg S)+r(T\wedge\neg S)\) by additivity. So
\(s_r(T,S)=1\) exactly when \(r(T\wedge\neg S)=0\).

(2)\(\Leftrightarrow\)(3). By \cref{def:p15-r43-extreme}, \(T\) is
extremely positive to \(S\) when \(P_r(T)>0\), \(P_r(S)<1\) and
\(P_r(S\mid T)=1\). The first two hold by hypothesis, and
\(P_r(S\mid T)=1\) means \(r(T\wedge S)=r(T)\), that is,
\(r(T\wedge\neg S)=0\).

(2)\(\Leftrightarrow\)(4). \(\neg S\) is extremely positive to \(\neg T\)
when \(P_r(\neg S)>0\), \(P_r(\neg T)<1\) and
\(P_r(\neg T\mid\neg S)=1\). The first two are \(r(S)<1\) and
\(r(T)>0\). The third says \(r(\neg T\wedge\neg S)=r(\neg S)\), which by
the additivity used for (1)\(\Leftrightarrow\)(2) is
\(r(T\wedge\neg S)=0\).

If \(T\entails S\), then \(T\wedge\neg S\) is contradictory and
\(r(T\wedge\neg S)=r(\bot)=0\).

Values. Assume (2). Then \(r(T)=r(T\wedge S)\le r(S)<1\). So
\(p_r(T,S)=r(T)/r(S)\) and
\(R_D(S\Yright T)_r=r(T)/r(S)-r(T)=r(T)(1-r(S))/r(S)\). Since
\(r(T)>0\), \(R_Q(S\Yright T)_r=p_r(T,S)/r(T)=1/r(S)\). Since
\(0<r(T)<1\), the likelihood ratio is defined: the likelihood under
\(T\) is \(P_r(S\mid T)=1\), and under \(\neg T\) it is
\[
 P_r(S\mid\neg T)=\frac{r(S\wedge\neg T)}{r(\neg T)}
 =\frac{r(S)-r(T)}{1-r(T)} ,
\]
using \(r(S)=r(S\wedge T)+r(S\wedge\neg T)\) and \(r(S\wedge T)=r(T)\).
If \(r(S)>r(T)\) this is positive and
\(R_L(S\Yright T)_r=(1-r(T))/(r(S)-r(T))\); if \(r(S)=r(T)\) it is zero
while the other likelihood is \(1\), and \cref{def:p15-r43-measures}
gives \(R_L=+\infty\). Finally \(s_r(T,S)-g(T)=1-(1-r(T))=r(T)\).

Monotonicity. Fix \(S\) and let \(t=r(T)\) range over \((0,r(S)]\).
\(t/r(S)\) and \(t(1-r(S))/r(S)\) are strictly increasing because
\(0<r(S)<1\). For \(t<r(S)\),
\[
 \frac{d}{dt}\,\frac{1-t}{r(S)-t}=\frac{1-r(S)}{(r(S)-t)^2}>0,
\]
and the value \(+\infty\) at \(t=r(S)\) exceeds every finite value. The
values \(s_r(T,S)=1\) and \(R_Q=1/r(S)\) do not depend on \(t\).
\end{proof}
\end{annexitem}

\begin{annexitem}{annex:p15-w02b-p15-025-not-popd}{thm:p15-w02b-p15-025-not-popd}{Power against POP-D}
\begin{proof}
The eight masses are \(P(H)\,P(E_1^{\pm}\mid H)\,P(E_2^{\pm}\mid H)\)
and the corresponding products under \(\neg H\), where \(E^{+}=E\) and
\(E^{-}=\neg E\). They are nonnegative and sum to
\(\tfrac12\cdot1\cdot1+\tfrac12\cdot1\cdot1=1\), so they define a
probability on the eight assignments; by
\cref{thm:p15-power-range-laws} it is a range measure on the singular
sentences.

For \(E_1\):
\(r(E_1)=\tfrac12\cdot\tfrac34+\tfrac12\cdot\tfrac14=\tfrac12\),
\(r(H\wedge E_1)=\tfrac38\), so \(p_r(H,E_1)=\tfrac34\). Also
\(r(\neg H\wedge\neg E_1)=\tfrac12\cdot\tfrac34=\tfrac38\) and
\(r(\neg E_1)=\tfrac12\), so \(s_r(H,E_1)=\tfrac34\) by
\cref{thm:p15-w02b-p15-025-common-model}.

For \(E_2\):
\(r(E_2)=\tfrac12\cdot\tfrac7{16}+\tfrac12\cdot\tfrac1{16}=\tfrac14\),
\(r(H\wedge E_2)=\tfrac7{32}\), so \(p_r(H,E_2)=\tfrac78\). Also
\(r(\neg H\wedge\neg E_2)=\tfrac12\cdot\tfrac{15}{16}=\tfrac{15}{32}\) and
\(r(\neg E_2)=\tfrac34\), so \(s_r(H,E_2)=\tfrac{15}{32}\cdot\tfrac43
=\tfrac58\).

With \(g(H)=\tfrac12\) the last column is \(\tfrac14\) and \(\tfrac18\).
As a check against \cref{thm:p15-w02b-p15-025-increment}:
\(O(\tfrac12)(\tfrac34-\tfrac12)=\tfrac14\) and
\(O(\tfrac14)(\tfrac78-\tfrac12)=\tfrac13\cdot\tfrac38=\tfrac18\).

The events \(H,E_1,E_2\) and their negations have probability
strictly between zero and one, so every measure in
\cref{def:p15-r43-measures} is in its defined domain. POP-D
(\cref{sec:p15-r43-axioms}) requires
\(R(E_2\Yright H)_P>R(E_1\Yright H)_P\) because
\(P(H\mid E_2)=\tfrac78>\tfrac34=P(H\mid E_1)\). Both
\(s_r(H,\cdot)\) and \(s_r(H,\cdot)-g(H)\) give the opposite strict
order. Ordinal equivalence requires every comparison to be preserved,
so neither function is ordinally equivalent on these fixed-target
comparisons to any measure satisfying POP-D.
\end{proof}
\end{annexitem}

\begin{annexitem}{annex:p15-w02b-p15-025-certain}{thm:p15-w02b-p15-025-certain}{Certain data}
\begin{proof}
Item~1 is the domain statement of
\cref{thm:p15-w02b-p15-025-common-model}.

Item~2. With weights \(w_v>0\) on \(\Omega\), \(r(S)=\sum_{v\in[S]}w_v\)
and \(\sum_{v\in\Omega}w_v=1\). Thus \(r(S)<1\) exactly when
\([S]\neq\Omega\), and \(r(S)>0\) exactly when \([S]\neq\varnothing\).
By \cref{thm:p15-power-normal-forms}, \(S\) is valid exactly when
\([S]=\Omega\) and contradictory exactly when \([S]=\varnothing\).

Item~3. Let \(r(S)>0\) and \(r_S(U)=r(U\wedge S)/r(S)\). Values lie in
\([0,1]\) because \(r(U\wedge S)\le r(S)\). If \(U\) is valid,
\(U\wedge S\equiv S\) and \(r_S(U)=1\). If \(U\wedge V\) is
contradictory, so is \((U\wedge S)\wedge(V\wedge S)\), and
\((U\vee V)\wedge S\equiv(U\wedge S)\vee(V\wedge S)\), so
\(r_S(U\vee V)=r_S(U)+r_S(V)\). Hence \(r_S\) satisfies
\cref{def:p15-power-range-measure}. \(r_S(S)=r(S)/r(S)=1\), so
\(s_{r_S}(T,S)\) is undefined for every \(T\) by item~1. For every
\(U\), \(\Delta_{r_S}(S,U)=r_S(S\wedge U)-r_S(S)\,r_S(U)
=r_S(U)-r_S(U)=0\), so \(S\) is irrelevant to \(U\) under \(P_{r_S}\).
\end{proof}
\end{annexitem}

\begin{annexitem}{annex:p15-w02b-p15-025-example}{thm:p15-w02b-p15-025-example}{The worked example}
\begin{proof}
List the assignments to \((P(a),Q(a))\) as \(11,10,01,00\), each of
weight \(1/4\) under \(r_1\).

Item~1. \(S=P(a)\wedge Q(a)\) holds only at \(11\), so \(r_1(S)=1/4\).
\(T^\flat=\neg P(a)\vee Q(a)\) holds at \(11,01,00\), so
\(r_1(T^\flat)=3/4\) and \(g_1(T^\flat)=1/4\). Because \(S\entails
T^\flat\), \(T^\flat\vee S\equiv T^\flat\), and
\(s_{r_1}(T^\flat,S)=g_1(T^\flat)/g_1(S)=(1/4)/(3/4)=1/3\). By
\cref{thm:p15-power-uniform-formula}, \(s_0(T,K)\) equals this value,
and \cref{thm:p15-power-source-examples} gives \(s_0(T,K)=1/3\)
directly. \(p_{r_1}(T^\flat,S)=r_1(T^\flat\wedge S)/r_1(S)=1\).
\(R_D(S\Yright T^\flat)_{r_1}=1-3/4=1/4\). For \(R_L\): the target has
probability \(3/4\in(0,1)\); \(P(S\mid T^\flat)=(1/4)/(3/4)=1/3>0\);
\(\neg T^\flat\) holds only at \(10\), where \(S\) is false, so
\(P(S\mid\neg T^\flat)=0\) and \(R_L=+\infty\) by
\cref{def:p15-r43-measures}. Since \(r_1(S)>0\), \(r_1(T^\flat)<1\) and
\(P(T^\flat\mid S)=1\), \(S\) is extremely positive to \(T^\flat\)
(\cref{def:p15-r43-extreme}). Finally \(O(1/4)=1/3\) and
\(1/3\cdot1/4=1/12=1/3-1/4\).

Item~2. \(S'=P(a)\vee Q(a)\) holds at \(11,10,01\), so
\(r_1(S')=3/4\). \(T_1\) holds only at \(11\), so \(r_1(T_1)=1/4\);
\(T_2\) holds at \(11,10\), so \(r_1(T_2)=1/2\). Both entail \(S'\),
and both have \(r_1(T_i)>0\) and \(0<r_1(S')<1\), so
\cref{thm:p15-w02b-p15-025-conclusive} applies:
\(s_{r_1}(T_i,S')=1\); \(p_{r_1}=r_1(T_i)/(3/4)\), giving \(1/3\) and
\(2/3\); \(R_D=r_1(T_i)\cdot(1/4)/(3/4)=r_1(T_i)/3\), giving \(1/12\)
and \(1/6\); \(R_Q=4/3\) for both; \(R_L=(1-r_1(T_i))/(3/4-r_1(T_i))\),
giving \((3/4)/(1/2)=3/2\) and \((1/2)/(1/4)=2\); and
\(s_{r_1}-g_1=r_1(T_i)\), giving \(1/4\) and \(1/2\).
\end{proof}
\end{annexitem}

The remaining items prove the results of
\cref{sec:p15-w02b-p15-025-ml,sec:p15-w02b-p15-025-variance,sec:p15-w02b-p15-025-confidence,sec:p15-w02b-p15-025-minimax}.
All sums are finite unless an integral is written.

\begin{annexitem}{annex:p15-w02b-p15-025-ml-cstar}{thm:p15-w02b-p15-025-ml-cstar}{The two frequency estimates}
\begin{proof}
Item~1. For \(s\ge1\),
\[
 \frac{s_M+w}{s+k}-\frac{s_M}{s}
 =\frac{s(s_M+w)-(s+k)s_M}{s(s+k)}
 =\frac{ws-ks_M}{s(s+k)} .
\]
The numerator is zero exactly when \(ws=ks_M\), that is, when
\(s_M/s=w/k\). As \(s_M\) runs over \(0,1,\ldots,s\), the numerator
\(ws-ks_M\) decreases from \(ws\) at \(s_M=0\) to \(ws-ks=-s(k-w)\) at
\(s_M=s\), so its absolute value is at most \(s\max(w,k-w)\), with
equality at \(s_M=0\) when \(w\ge k-w\) and at \(s_M=s\) when
\(w\le k-w\). Dividing by \(s(s+k)>0\) gives the bound and its cases of
equality.

Item~2. In the Bernoulli model of \cref{def:p15-fe-binomial} with
\(n=s\) trials and \(x=s_M\) successes, \cref{thm:p15-fe-binomial}
gives the unique maximizer \(\widehat p=x/n\) on \([0,1]\); this is
\(0\) when \(s_M=0\) and \(1\) when \(s_M=s\). On Carnap's side,
\(s_M=0\) gives \(c^*=w/(s+k)\), positive because \(w>0\); and
\(s_M=s\) gives \(c^*=(s+w)/(s+k)\), less than \(1\) because \(w<k\).

Item~3. With no observations the likelihood of the empty sample is the
empty product, equal to \(1\) for every \(p\in[0,1]\), so every
parameter value maximizes it and no value is the unique maximizer. The
formula of \cref{thm:p15-carnap-predictive} with \(s=s_M=0\) gives
\(c^*=w/k\).
\end{proof}
\end{annexitem}

\begin{annexitem}{annex:p15-w02b-p15-025-variance}{thm:p15-w02b-p15-025-variance}{Three variance multipliers}
\begin{proof}
For any reals \(x_1,\ldots,x_n\) with mean \(\bar x\) and any \(m\),
expanding \((x_i-m)^2=((x_i-\bar x)+(\bar x-m))^2\) and using
\(\sum_i(x_i-\bar x)=0\) gives
\[
 \sum_{i=1}^n(x_i-m)^2=\sum_{i=1}^n(x_i-\bar x)^2+n(\bar x-m)^2 .
 \tag{\(*\)}
\]

Item~1. By \((*)\) with \(m=\mu\),
\(V=\sum_i(X_i-\mu)^2-n(\bar X-\mu)^2\). Each \(\mathbb E(X_i-\mu)^2\)
is \(\sigma^2\). By independence the covariances
\(\mathbb E[(X_i-\mu)(X_j-\mu)]\), \(i\ne j\), vanish, so
\(\mathbb E(\bar X-\mu)^2=n^{-2}\sum_i\mathbb E(X_i-\mu)^2=\sigma^2/n\).
Hence \(\mathbb EV=n\sigma^2-\sigma^2=\nu\sigma^2\), and dividing by
\(n\) and by \(\nu\) gives the two expectations.

Item~2. For observed \(x\) and parameters \(\mu\in\mathbb R\),
\(t=\sigma^2>0\), the likelihood is
\[
 \begin{aligned}
 L(\mu,t)&=(2\pi t)^{-n/2}\exp\Bigl(-\frac1{2t}\sum_i(x_i-\mu)^2\Bigr)\\
 &=(2\pi t)^{-n/2}\exp\Bigl(-\frac{v+n(\bar x-\mu)^2}{2t}\Bigr)
 \end{aligned}
\]
by \((*)\). For fixed \(t\), it is largest exactly at \(\mu=\bar x\),
and strictly smaller elsewhere. Put
\(\ell(t)=\log L(\bar x,t)=-\tfrac n2\log(2\pi)-\tfrac n2\log t-v/(2t)\).
Then
\[
 \ell'(t)=-\frac n{2t}+\frac v{2t^2}=\frac{v-nt}{2t^2}.
\]
If the observations are not all equal then \(v>0\), \(\ell'\) is
positive on \((0,v/n)\) and negative on \((v/n,\infty)\), so \(t=v/n\)
is the unique maximizer of \(\ell\). For
\((\mu,t)\ne(\bar x,v/n)\),
\(L(\mu,t)\le L(\bar x,t)\le L(\bar x,v/n)\), and one inequality is
strict: the first if \(\mu\ne\bar x\), the second if \(\mu=\bar x\) and
\(t\ne v/n\). If all observations are equal, then \(v=0\) and
\(L(\bar x,t)=(2\pi t)^{-n/2}\) tends to \(+\infty\) as \(t\to0\), so
the likelihood is unbounded and has no maximizer.

Item~3. By \cref{thm:p15-student-orthogonal}, \(V/\sigma^2=W=\sum_{j=2}^nY_j^2\)
with \(Y_2,\ldots,Y_n\) independent standard normal. Let \(\varphi\) be
the standard normal density; \(\varphi'(y)=-y\varphi(y)\), and
\(y^m\varphi(y)\to0\) as \(|y|\to\infty\) for every \(m\). Integration
by parts gives
\[
 \begin{gathered}
 \mathbb EY^2=\int y\cdot y\varphi(y)\,dy=\int\varphi(y)\,dy=1,\\
 \mathbb EY^4=\int y^3\cdot y\varphi(y)\,dy=3\int y^2\varphi(y)\,dy=3,
 \end{gathered}
\]
so \(\operatorname{Var}(Y^2)=3-1=2\). Hence \(\mathbb EW=\nu\) and, by
independence, \(\operatorname{Var}W=2\nu\) and
\(\mathbb EW^2=2\nu+\nu^2=\nu(\nu+2)\). For \(a>0\),
\[
 \mathbb E(aV-\sigma^2)^2=\sigma^4\,\mathbb E(aW-1)^2
 =\sigma^4\bigl(a^2\nu(\nu+2)-2a\nu+1\bigr).
\]
This is a quadratic in \(a\) with positive leading coefficient, so its
unique minimizer is \(a=\nu/(\nu(\nu+2))=1/(\nu+2)=1/(n+1)\). The
bracket equals \(1-\nu/(\nu+2)=2/(n+1)\) at \(a=1/(\nu+2)\);
\((\nu+2)/\nu-1=2/\nu=2/(n-1)\) at \(a=1/\nu\); and
\[
 \frac{\nu(\nu+2)-2\nu(\nu+1)+(\nu+1)^2}{(\nu+1)^2}
 =\frac{2\nu+1}{(\nu+1)^2}=\frac{2n-1}{n^2}
\]
at \(a=1/(\nu+1)=1/n\). Finally \(2n^2<(2n-1)(n+1)=2n^2+n-1\) because
\(n>1\), and \((2n-1)(n-1)=2n^2-3n+1<2n^2\), which give the two strict
inequalities.
\end{proof}
\end{annexitem}

\begin{annexitem}{annex:p15-w02b-p15-025-coverage}{thm:p15-w02b-p15-025-coverage}{Coverage, direct inference and averaged confirmation}
\begin{proof}
Write \(Q=Q_\pi\) and \(A=\{(\theta,x):\theta\in C(x)\}\).

Item~1. \(Q(\theta=\theta_0)=\pi(\theta_0)\sum_xP_{\theta_0}(x)=\pi(\theta_0)>0\),
and
\(Q(A\cap\{\theta=\theta_0\})=\sum_{x:\,\theta_0\in C(x)}\pi(\theta_0)P_{\theta_0}(x)
=\pi(\theta_0)\gamma_C(\theta_0)\). Their ratio is \(\gamma_C(\theta_0)\).

Item~2. For \(Q(x)>0\),
\(Q(x)\,Q(C(x)\mid x)=Q(A\cap(\Theta\times\{x\}))=\sum_{\theta\in C(x)}\pi(\theta)P_\theta(x)\).
If \(Q(x)=0\), every term \(\pi(\theta)P_\theta(x)\) is zero, so adding
those \(x\) changes nothing. Exchanging the two finite sums,
\[
 \begin{aligned}
 \sum_{x:\,Q(x)>0}Q(x)\,Q(C(x)\mid x)
 &=\sum_x\sum_{\theta\in C(x)}\pi(\theta)P_\theta(x)\\
 &=\sum_\theta\pi(\theta)\sum_{x:\,\theta\in C(x)}P_\theta(x)
 =\sum_\theta\pi(\theta)\gamma_C(\theta).
 \end{aligned}
\]
Since \(\sum_{x:\,Q(x)>0}Q(x)=1\), the left side is the
\(Q\)-expectation of the confirmation of the reported set. If
\(\gamma_C\ge1-\alpha\), the right side is at least
\((1-\alpha)\sum_\theta\pi(\theta)=1-\alpha\), with equality when
\(\gamma_C=1-\alpha\) everywhere.

Item~3. \(0\in C(x)\) only for \(x=1\), so
\(\gamma_C(0)=P_0(1)=9/10\); \(1\in C(x)\) only for \(x=3\), so
\(\gamma_C(1)=P_1(3)=9/10\). Let \(\pi(0),\pi(1)>0\). Then
\(Q(1)=\tfrac9{10}\pi(0)>0\) and all of it lies on \(\theta=0\in C(1)\),
so the confirmation of \(C(1)\) given \(x=1\) is \(1\); symmetrically
\(Q(3)=\tfrac9{10}\pi(1)>0\) lies on \(\theta=1\in C(3)\), giving \(1\).
\(Q(2)=\tfrac1{10}(\pi(0)+\pi(1))=\tfrac1{10}>0\) and \(C(2)\) is empty,
giving \(0\). As a check on item~2, the average is
\(\tfrac9{10}\pi(0)+\tfrac9{10}\pi(1)=\tfrac9{10}\).
\end{proof}
\end{annexitem}

\begin{annexitem}{annex:p15-w02b-p15-025-minimax}{thm:p15-w02b-p15-025-minimax}{Minimax and weighted loss}
\begin{proof}
The expected loss of \(\delta_t\) is \(4(1-t)\) at \(\theta_1\) and
\(10t+4(1-t)=4+6t\) at \(\theta_2\). For \(t\in[0,1]\),
\(4+6t\ge4-4t\), so the maximum is \(4+6t\), which is smallest exactly
at \(t=0\), with value \(4\). This proves item~1.

With confirmation \(1-q\) for \(\theta_1\) and \(q\) for \(\theta_2\),
the weighted loss of \(\delta_t\) is
\((1-q)(4-4t)+q(4+6t)=4+t(10q-4)\). This is linear in \(t\) with slope
\(10q-4\): negative for \(q<2/5\), so \(t=1\) (the action \(a_1\)) is
the unique minimizer; positive for \(q>2/5\), so \(t=0\) (the action
\(a_2\)) is; zero for \(q=2/5\), so every \(t\) is. Comparing with
item~1, the minimax choice \(a_2\) differs from the weighted choice
\(a_1\) for \(q<2/5\) and coincides with it for \(q>2/5\).
\end{proof}
\end{annexitem}
